\scanpage{91}
\mspage{45}
$\mathcal E_U$ sur $\mathcal E_S$ \unsure{ou|où} $\unsure{=}$ de $\mathcal E_{\widetilde U}$ sur $\unsure{\mathcal E_{\widetilde\sigma}|\mathcal E_{\widetilde S}}$, définit néanmoins, \deleted{\unsure{le}} le sous-groupe $\mathcal E_{\widetilde U}=N_{s'}$, de $\mathcal E_K$ (qui est extension d'un $\mathcal E_k$ par un facteur trivial \unsure{kummérien|kummerien} $T\simeq I_{s'}$, où $k$ est un corps de degré de \unsure{transcendance} $\unsure{\deg\operatorname{tr}K-1|\deg\operatorname{tr}n-1}$, lorsque $S$ est de type fini sur un corps de $\unsure{N}$, disons $\mathbf Q$) \uline{presque} un relèvement --- seul le facteur $I_{s'}\simeq T$ y \unsure{obstrue} encore : en \unsure{tous|tout} cas $N_{s'}\simeq\mathcal E_{\widetilde U}$ est plongé dans un $N_{\mathcal E_{X_s}}(L_i)$, et toute section de \unsure{cette} extension de $\unsure{E_{k(s)}|\mathcal E_{k(s)}}$ par $T$ donne donc une section de $\mathcal E_{X_s}$ sur $\mathcal E_{k(s)}$ --- mais de \uline{deuxième} espèce.

Une autre façon de dire ces choses est la suivante. Une section rationnelle $f$ de $\frac{U}{S}$ sur $S$, de domaine de définition l'ouvert $S'$ de $S$, \deleted{\unsure{peut}} peut s'interpréter en termes de la famille des pts $f(s)\in U_s$, où $s$ parcourt l'ens des pts de $S$
\margin{90}
